% =============================================================================
% Architecture Figure (includable version)
% Usage:  % =============================================================================
% Architecture Figure (includable version)
% Usage:  % =============================================================================
% Architecture Figure (includable version)
% Usage:  % =============================================================================
% Architecture Figure (includable version)
% Usage:  \input{architecture_figure.tex}
% Place inside your paper .tex file — requires tikz + xcolor already loaded
% Packages needed in preamble:
%   \usepackage{tikz}
%   \usetikzlibrary{arrows.meta,positioning,shapes.geometric,calc,
%                    fit,backgrounds,decorations.pathreplacing,patterns}
% =============================================================================

\begin{figure*}[t]
\centering
\resizebox{\textwidth}{!}{%
\input{architecture_diagram_body.tex}%
}
\caption{%
  Architecture of the proposed Generative VQA model with Mixture-of-Experts (MOE) fusion.
  Frozen CLIP ViT-B/32 and PhoBERT-base encode the image and Vietnamese question, respectively.
  Their projected embeddings are concatenated and processed through two pre-norm transformer layers,
  then routed via a Noisy Top-2 gating mechanism to $K{=}4$ heterogeneous specialised experts
  (Vision, Text, Multimodal, Specialized).
  Only $k{=}2$ experts are activated per token (50\% sparsity),
  maintaining computational efficiency while enabling adaptive multi-modal reasoning.
  The weighted expert outputs are decoded by a 6-layer autoregressive transformer with weight-tied embeddings.
  Dashed blue arrows denote sparse dispatch; dashed red arrows denote auxiliary loss paths.
  Gradients flow through the fusion, MOE, and decoder modules (dotted green);
  backbone encoders remain frozen.%
}
\label{fig:architecture}
\end{figure*}

% Place inside your paper .tex file — requires tikz + xcolor already loaded
% Packages needed in preamble:
%   \usepackage{tikz}
%   \usetikzlibrary{arrows.meta,positioning,shapes.geometric,calc,
%                    fit,backgrounds,decorations.pathreplacing,patterns}
% =============================================================================

\begin{figure*}[t]
\centering
\resizebox{\textwidth}{!}{%
% =============================================================================
% Architecture Diagram Body  (no preamble — for \input inside a figure)
% Requires in main document preamble:
%   \usepackage{tikz,amsmath,amssymb,xcolor}
%   \usetikzlibrary{arrows.meta,positioning,shapes.geometric,calc,
%                    fit,backgrounds,decorations.pathreplacing,patterns}
% =============================================================================

% ── Colours ──────────────────────────────────────────────────────────────────
\definecolor{visBlue}{HTML}{2E6DAD}
\definecolor{visBg}{HTML}{D6E6F5}
\definecolor{txtGreen}{HTML}{237A45}
\definecolor{txtBg}{HTML}{D4EDDE}
\definecolor{moeRed}{HTML}{C0392B}
\definecolor{moeBg}{HTML}{FDECEC}
\definecolor{moeBorder}{HTML}{D45A5A}
\definecolor{routerOrange}{HTML}{D98B2B}
\definecolor{routerBg}{HTML}{FDF3E2}
\definecolor{expPurple}{HTML}{7B5EA7}
\definecolor{expBg}{HTML}{EDE6F6}
\definecolor{decPurple}{HTML}{6A4C93}
\definecolor{decBg}{HTML}{E8E0F0}
\definecolor{outGray}{HTML}{4A4A4A}
\definecolor{outBg}{HTML}{F0F0F0}
\definecolor{arrGray}{HTML}{555555}
\definecolor{dimGray}{HTML}{888888}
\definecolor{lossRed}{HTML}{B03A2E}
\definecolor{gradGreen}{HTML}{2E7D32}
\definecolor{sparseBlue}{HTML}{1565C0}
\definecolor{frozenGray}{HTML}{95A5A6}
\definecolor{inputBlue}{HTML}{4A90D9}
\definecolor{inputBg}{HTML}{EBF2FA}

% ── Styles ───────────────────────────────────────────────────────────────────
\tikzset{
  blk/.style={draw,rounded corners=3pt,minimum height=0.95cm,
              font=\small\sffamily,align=center,thick,inner sep=6pt},
  inblk/.style ={blk,fill=inputBg,draw=inputBlue,minimum width=3.2cm},
  vblk/.style  ={blk,fill=visBg,draw=visBlue,minimum width=3.6cm},
  tblk/.style  ={blk,fill=txtBg,draw=txtGreen,minimum width=3.6cm},
  rblk/.style  ={blk,fill=routerBg,draw=routerOrange,minimum width=5.0cm},
  eblk/.style  ={blk,fill=expBg,draw=expPurple,
                  minimum width=2.4cm,minimum height=1.15cm,
                  font=\footnotesize\sffamily},
  dblk/.style  ={blk,fill=decBg,draw=decPurple,minimum width=6.0cm},
  oblk/.style  ={blk,fill=outBg,draw=outGray,minimum width=3.8cm},
  op/.style     ={draw,rounded corners=2pt,minimum height=0.55cm,
                  font=\scriptsize\sffamily,align=center,thin,
                  inner sep=4pt,fill=white},
  dim/.style    ={font=\scriptsize\ttfamily,text=dimGray},
  frozen/.style ={font=\tiny\sffamily\itshape,text=frozenGray},
  stg/.style    ={font=\footnotesize\sffamily\bfseries,text=dimGray},
  arr/.style    ={-{Stealth[length=5pt,width=4pt]},thick,arrGray},
  sparr/.style  ={-{Stealth[length=5pt,width=4pt]},thick,sparseBlue,dashed},
  auxarr/.style ={-{Stealth[length=4pt,width=3pt]},semithick,lossRed,
                  densely dashed},
  gradarr/.style={-{Stealth[length=4pt,width=3pt]},semithick,gradGreen,
                  densely dotted},
}

\begin{tikzpicture}[node distance=0.8cm and 1.5cm]

% ═══════════════════════════════════════════════════════════════════════════════
%  Y-COORDINATE MAP (generous spacing)
%
%   0.0   Inputs
%  -1.8   Patch / Token Embedding
%  -3.4   Transformer Encoder backbones
%  -5.0   Linear projections
%  -6.8   Concatenation
%  -8.8   Pre-Norm Transformer (fusion)
% -10.8   Noisy Top-2 Router
% -13.2   Expert blocks
% -15.4   Weighted Aggregation
% -17.6   Decoder
% -19.4   Output Projection
% -20.4   Softmax
% -21.6   Answer
% ═══════════════════════════════════════════════════════════════════════════════


% ═══════════════════  STAGE 1 — INPUTS  ══════════════════════════════════════
\node[stg] at (-8.0,0) {\rotatebox{90}{INPUT}};

\node[inblk] (imgIn) at (-3.2,0) {\textbf{Image}~$I$};
\node[dim,below=4pt of imgIn] {$224{\times}224{\times}3$};

\node[inblk] (qIn) at (3.2,0) {\textbf{Question}~$Q$};
\node[dim,below=4pt of qIn] {$L$ tokens (Vietnamese)};


% ═══════════════════  STAGE 2/3 — ENCODERS  ═════════════════════════════════
\node[stg] at (-8.0,-3.4) {\rotatebox{90}{ENCODERS}};

% --- Visual column (left) ---
\node[vblk] (patchEmb) at (-3.2,-1.8)
  {Patch Embedding\\[-2pt]{\scriptsize $32{\times}32$ patches + \texttt{[CLS]}}};
\node[dim,right=8pt of patchEmb] {$50{\times}d_v$};

\node[vblk] (vitEnc) at (-3.2,-3.4)
  {\textbf{CLIP ViT-B/32}\\[-2pt]{\scriptsize Transformer Encoder}};
\node[frozen,right=8pt of vitEnc] {\textit{frozen}};

\node[op,draw=visBlue,fill=visBg!40] (visProj) at (-3.2,-5.0)
  {Linear~$W_v$};
\node[dim,right=8pt of visProj] {$50{\times}D$};

% --- Text column (right) ---
\node[tblk] (tokEmb) at (3.2,-1.8)
  {Token Embedding\\[-2pt]{\scriptsize + Position Encoding}};
\node[dim,left=8pt of tokEmb] {$L{\times}d_t$};

\node[tblk] (phoBERT) at (3.2,-3.4)
  {\textbf{PhoBERT-base}\\[-2pt]{\scriptsize Transformer Encoder}};
\node[frozen,left=8pt of phoBERT] {\textit{frozen}};

\node[op,draw=txtGreen,fill=txtBg!40] (txtProj) at (3.2,-5.0)
  {Linear~$W_t$};
\node[dim,left=8pt of txtProj] {$L{\times}D$};

% Encoder arrows
\draw[arr] (imgIn)    -- (patchEmb);
\draw[arr] (patchEmb) -- (vitEnc);
\draw[arr] (vitEnc)   -- (visProj);
\draw[arr] (qIn)      -- (tokEmb);
\draw[arr] (tokEmb)   -- (phoBERT);
\draw[arr] (phoBERT)  -- (txtProj);


% ═══════════════════  CONCATENATION  ═════════════════════════════════════════
\node[op,minimum width=5.2cm,draw=arrGray] (concat) at (0,-6.8)
  {Concatenation~$[\hat{V};\;\hat{T}]$};
\node[dim,below=4pt of concat] (concatDim) {$(N_v{+}L){\times}D$};

\draw[arr] (visProj.south) -- ++(0,-0.5) -| ([xshift=-1.2cm]concat.north);
\draw[arr] (txtProj.south) -- ++(0,-0.5) -| ([xshift= 1.2cm]concat.north);


% ═══════════════════  STAGE 4 — FUSION WITH MOE  ════════════════════════════
\node[stg] at (-8.0,-12.0) {\rotatebox{90}{FUSION~(MOE)}};

% (a) Pre-norm Transformer attention
\node[blk,fill=moeBg!60,draw=moeRed!70,minimum width=6.0cm]
  (preAttn) at (0,-8.8)
  {\textbf{Pre-Norm Transformer}~${\times}\,L_f{=}2$\\[-2pt]
   {\scriptsize MHSA~($H{=}8$)\;+\;FFN~($d_{\mathrm{ff}}{=}2048$)\;+\;LN}};
\node[dim,right=10pt of preAttn] {$(N_v{+}L){\times}D$};

% (b) Noisy Top-k Router
\node[rblk] (router) at (0,-10.8)
  {\textbf{Noisy Top-2 Router}\\[-2pt]
   {\scriptsize $\tilde{\ell}=W_r^\top x
     +\mathrm{Softplus}(W_n^\top x)\!\cdot\!\epsilon$\;,\;\;
     $\epsilon{\sim}\mathcal{N}(0,\sigma^2 I)$}\\[-2pt]
   {\scriptsize Softmax\;$\rightarrow$\;Top-2 selection per token}};
\node[dim,right=10pt of router,text=sparseBlue]
  {\textit{sparse:\;$k{=}2$ of $K{=}4$}};

% (c) Specialised Experts — wider horizontal spread
\node[eblk] (eVis)  at (-4.5,-13.2)
  {$E_{\mathrm{vis}}$\\[-1pt]\textbf{Vision}\\[-2pt]
   {\tiny Spatial MHSA}\\[-2pt]{\tiny + FFN}};
\node[eblk] (eTxt)  at (-1.5,-13.2)
  {$E_{\mathrm{txt}}$\\[-1pt]\textbf{Text}\\[-2pt]
   {\tiny Context MHSA}\\[-2pt]{\tiny + Masked FFN}};
\node[eblk] (eMM)   at ( 1.5,-13.2)
  {$E_{\mathrm{mm}}$\\[-1pt]\textbf{Multimodal}\\[-2pt]
   {\tiny Cross-Attn}\\[-2pt]{\tiny + Gate~$\gamma$}};
\node[eblk] (eSpec) at ( 4.5,-13.2)
  {$E_{\mathrm{spec}}$\\[-1pt]\textbf{Specialized}\\[-2pt]
   {\tiny Seg\,/\,Det\,/}\\[-2pt]{\tiny OCR\,/\,Scene}};

% (d) Weighted aggregation
\node[op,minimum width=7.0cm,draw=moeRed!70,fill=moeBg!50]
  (agg) at (0,-15.4)
  {Weighted Aggregation\;+\;LayerNorm};


% Load-balance loss (side branch — pushed further right)
\node[blk,fill=white,draw=lossRed,minimum width=2.6cm,
      font=\scriptsize\sffamily] (lbLoss) at (8.0,-10.8)
  {$\mathcal{L}_{\mathrm{lb}}
     =K\!\displaystyle\sum f_i\bar{p}_i$\\[0pt]
   {\tiny $\lambda_{\mathrm{lb}}{=}0.01$}};
\node[dim,above=5pt of lbLoss,text=lossRed]
  {\textit{Load-Balance Loss}};

% ── MOE highlighted background frame ────────────────────────────────────────
\begin{scope}[on background layer]
  \node[draw=moeBorder,fill=moeBg!20,rounded corners=6pt,
        very thick,inner sep=14pt,
        fit=(preAttn)(router)(eVis)(eSpec)(agg)(lbLoss)]
    (moeFrame) {};
  \node[font=\small\sffamily\bfseries,text=moeRed,
        anchor=north west]
    at ([xshift=6pt,yshift=-1pt]moeFrame.north west)
    {Cross-Modal Fusion with Mixture-of-Experts};
\end{scope}

% ── Arrows inside MOE ───────────────────────────────────────────────────────
\draw[arr] (concat) -- (preAttn);
\draw[arr] (preAttn) -- (router);

% Router → experts  (sparse dashed blue)
\draw[sparr] (router.south) -- ++(0,-0.5) -| (eVis.north);
\draw[sparr] (router.south) -- ++(0,-0.5) -| (eTxt.north);
\draw[sparr] (router.south) -- ++(0,-0.5) -| (eMM.north);
\draw[sparr] (router.south) -- ++(0,-0.5) -| (eSpec.north);

% Experts → aggregation
\draw[arr] (eVis.south)  -- ++(0,-0.45) -| ([xshift=-2.1cm]agg.north);
\draw[arr] (eTxt.south)  -- ++(0,-0.50) -| ([xshift=-0.7cm]agg.north);
\draw[arr] (eMM.south)   -- ++(0,-0.50) -| ([xshift= 0.7cm]agg.north);
\draw[arr] (eSpec.south) -- ++(0,-0.45) -| ([xshift= 2.1cm]agg.north);

% Router → load-balance loss
\draw[auxarr] (router.east) -- (lbLoss.west);


% ═══════════════════  STAGE 5 — DECODER  ════════════════════════════════════
\node[stg] at (-8.0,-18) {\rotatebox{90}{DECODER}};

\node[dblk] (decoder) at (0,-18)
  {\textbf{Autoregressive Transformer Decoder}~${\times}\,L_d{=}6$\\[-2pt]
   {\scriptsize Causal Self-Attn\;$\to$\;Cross-Attn\;$\to$\;FFN}\\[-2pt]
   {\scriptsize $H_d{=}8$ heads,\;\; $d_{\mathrm{ff}}{=}2048$,\;\; Pre-Norm}};
\node[dim,right=10pt of decoder]
  {\parbox{2.4cm}{\centering\scriptsize encoder memory\\[1pt]
   $(N_v{+}L){\times}D$}};

% Teacher-forcing input (left, with more clearance)
\node[op,draw=decPurple,fill=decBg!40,minimum width=2.0cm]
  (tfEmb) at (-5.8,-18)
  {\parbox{1.8cm}{\centering Embed.~$W_e$\\[1pt]{\tiny + Pos Enc}\\[1pt]
   {\tiny (teacher forcing)}}};
\node[dim,above=1pt of tfEmb] {$a_{<t}$};

\draw[arr] (agg) -- (decoder);
\draw[arr] (tfEmb) -- (decoder);


% ═══════════════════  STAGE 6 — OUTPUT  ═════════════════════════════════════
\node[stg] at (-8.0,-20.6) {\rotatebox{90}{OUTPUT}};

\node[op,draw=outGray,minimum width=4.5cm] (outProj) at (0,-19.4)
  {Output Projection~$W_e^\top$\;\;{\scriptsize (weight-tied)}};
\node[dim,right=10pt of outProj] {$T{\times}|\mathcal{V}|$};

\node[op,draw=outGray,minimum width=2.8cm] (softmax) at (0,-20.4)
  {Softmax};

\node[oblk] (answer) at (0,-21.6)
  {\textbf{Generated Answer}\\[-2pt]
   {\scriptsize Open-ended Vietnamese text}};

\draw[arr] (decoder) -- (outProj);
\draw[arr] (outProj) -- (softmax);
\draw[arr] (softmax) -- (answer);


% ═══════════════════  TRAINING LOSS  ═════════════════════════════════════════
\node[blk,fill=white,draw=lossRed,minimum width=3.4cm,
      font=\scriptsize\sffamily] (totalLoss) at (8.0,-18.0)
  {\textbf{Training Loss}\\[-2pt]
   $\mathcal{L}=\mathcal{L}_{\mathrm{ce}}
     +\lambda\,\mathcal{L}_{\mathrm{lb}}$\\[-2pt]
   {\tiny label smoothing~$\epsilon{=}0.1$}};

\draw[auxarr] (softmax.east) -| (totalLoss.south);
\draw[auxarr] (lbLoss.south) |- ([yshift=4pt]totalLoss.north);


% ═══════════════════  GRADIENT FLOW  ════════════════════════════════════════
\draw[gradarr,very thick]
  ([xshift=8.3cm]answer.east) -- ++(0,20.0)
  node[midway,right=4pt,font=\scriptsize\sffamily,
       text=gradGreen,rotate=90,anchor=south]
  {$\nabla$~Gradient flow (backpropagation)};


% ═══════════════════  FROZEN BACKBONE NOTE  ═════════════════════════════════
\node[font=\tiny\sffamily,text=frozenGray,align=center,
      fill=white,draw=frozenGray,rounded corners=2pt,inner sep=3pt]
  (frozenNote) at (-6.2,-3.4)
  {\textsf{\textbf{$\times$}}~Frozen backbones:\\no gradient update};

\draw[frozenGray,thin,densely dotted]
  (frozenNote.east) -- ++(0.2,0) |- (vitEnc.west);
\draw[frozenGray,thin,densely dotted]
  (frozenNote.east) -- ++(0.2,0) |- (phoBERT.west);


% ═══════════════════  LEGEND (top-right)  ════════════════════════════════════
\node[anchor=north west,font=\scriptsize\sffamily,text=dimGray,
      align=left,fill=white,draw=dimGray!50,
      rounded corners=3pt,inner sep=6pt]
  at (10.5, 0.5)
  {\textbf{Legend}\\[4pt]
   \tikz[baseline=-0.5ex]\draw[arr](0,0)--(0.6,0);
     ~Dense data flow\\[2pt]
   \tikz[baseline=-0.5ex]\draw[sparr](0,0)--(0.6,0);
     ~Sparse dispatch\\[2pt]
   \tikz[baseline=-0.5ex]\draw[auxarr](0,0)--(0.6,0);
     ~Auxiliary loss\\[2pt]
   \tikz[baseline=-0.5ex]\draw[gradarr](0,0)--(0.6,0);
     ~Gradient flow\\[5pt]
   \tikz\fill[visBg](0,0) rectangle (0.4,0.25);
     ~Visual pathway\\[2pt]
   \tikz\fill[txtBg](0,0) rectangle (0.4,0.25);
     ~Text pathway\\[2pt]
   \tikz\fill[moeBg](0,0) rectangle (0.4,0.25);
     ~MOE fusion (core)\\[2pt]
   \tikz\fill[decBg](0,0) rectangle (0.4,0.25);
     ~Decoder};


% ═══════════════════  COMPUTATIONAL PROPERTIES (right)  ═════════════════════
\node[anchor=north west,font=\scriptsize\sffamily,text=dimGray,
      align=left,fill=white,draw=dimGray!50,
      rounded corners=3pt,inner sep=6pt]
  at (10.5,-5.0)
  {\textbf{Computational Properties}\\[4pt]
   \textcolor{sparseBlue}{\rule{8pt}{2pt}}~%
     Sparse activation: Top-2 of $K{=}4$\\[2pt]
   \textcolor{sparseBlue}{\rule{8pt}{2pt}}~%
     Sparsity ratio: $k/K=50\%$ per token\\[2pt]
   Capacity factor: $\alpha=1.25$\\[2pt]
   $D=d_v=d_t=768$\\[2pt]
   $|\mathcal{V}|=64{,}000$ (PhoBERT vocab)\\[2pt]
   Weight tying: $W_e$ shared in\\
   \quad decoder embed.\ \& output proj.\\[4pt]
   \textcolor{gradGreen}{\rule[1pt]{8pt}{1.5pt}}~%
     Gradients flow through\\
   \quad fusion, MOE, decoder\\[2pt]
   \textcolor{frozenGray}{\textbf{$\times$}}~%
     No gradient to frozen encoders\\[4pt]
   \textit{Dense fusion:}\;$\mathcal{O}(S^2 D)$\\[1pt]
   \textit{Sparse MOE:}\;$\mathcal{O}(S\!\cdot\!k\!\cdot\!D)$};

\end{tikzpicture}
%
}
\caption{%
  Architecture of the proposed Generative VQA model with Mixture-of-Experts (MOE) fusion.
  Frozen CLIP ViT-B/32 and PhoBERT-base encode the image and Vietnamese question, respectively.
  Their projected embeddings are concatenated and processed through two pre-norm transformer layers,
  then routed via a Noisy Top-2 gating mechanism to $K{=}4$ heterogeneous specialised experts
  (Vision, Text, Multimodal, Specialized).
  Only $k{=}2$ experts are activated per token (50\% sparsity),
  maintaining computational efficiency while enabling adaptive multi-modal reasoning.
  The weighted expert outputs are decoded by a 6-layer autoregressive transformer with weight-tied embeddings.
  Dashed blue arrows denote sparse dispatch; dashed red arrows denote auxiliary loss paths.
  Gradients flow through the fusion, MOE, and decoder modules (dotted green);
  backbone encoders remain frozen.%
}
\label{fig:architecture}
\end{figure*}

% Place inside your paper .tex file — requires tikz + xcolor already loaded
% Packages needed in preamble:
%   \usepackage{tikz}
%   \usetikzlibrary{arrows.meta,positioning,shapes.geometric,calc,
%                    fit,backgrounds,decorations.pathreplacing,patterns}
% =============================================================================

\begin{figure*}[t]
\centering
\resizebox{\textwidth}{!}{%
% =============================================================================
% Architecture Diagram Body  (no preamble — for \input inside a figure)
% Requires in main document preamble:
%   \usepackage{tikz,amsmath,amssymb,xcolor}
%   \usetikzlibrary{arrows.meta,positioning,shapes.geometric,calc,
%                    fit,backgrounds,decorations.pathreplacing,patterns}
% =============================================================================

% ── Colours ──────────────────────────────────────────────────────────────────
\definecolor{visBlue}{HTML}{2E6DAD}
\definecolor{visBg}{HTML}{D6E6F5}
\definecolor{txtGreen}{HTML}{237A45}
\definecolor{txtBg}{HTML}{D4EDDE}
\definecolor{moeRed}{HTML}{C0392B}
\definecolor{moeBg}{HTML}{FDECEC}
\definecolor{moeBorder}{HTML}{D45A5A}
\definecolor{routerOrange}{HTML}{D98B2B}
\definecolor{routerBg}{HTML}{FDF3E2}
\definecolor{expPurple}{HTML}{7B5EA7}
\definecolor{expBg}{HTML}{EDE6F6}
\definecolor{decPurple}{HTML}{6A4C93}
\definecolor{decBg}{HTML}{E8E0F0}
\definecolor{outGray}{HTML}{4A4A4A}
\definecolor{outBg}{HTML}{F0F0F0}
\definecolor{arrGray}{HTML}{555555}
\definecolor{dimGray}{HTML}{888888}
\definecolor{lossRed}{HTML}{B03A2E}
\definecolor{gradGreen}{HTML}{2E7D32}
\definecolor{sparseBlue}{HTML}{1565C0}
\definecolor{frozenGray}{HTML}{95A5A6}
\definecolor{inputBlue}{HTML}{4A90D9}
\definecolor{inputBg}{HTML}{EBF2FA}

% ── Styles ───────────────────────────────────────────────────────────────────
\tikzset{
  blk/.style={draw,rounded corners=3pt,minimum height=0.95cm,
              font=\small\sffamily,align=center,thick,inner sep=6pt},
  inblk/.style ={blk,fill=inputBg,draw=inputBlue,minimum width=3.2cm},
  vblk/.style  ={blk,fill=visBg,draw=visBlue,minimum width=3.6cm},
  tblk/.style  ={blk,fill=txtBg,draw=txtGreen,minimum width=3.6cm},
  rblk/.style  ={blk,fill=routerBg,draw=routerOrange,minimum width=5.0cm},
  eblk/.style  ={blk,fill=expBg,draw=expPurple,
                  minimum width=2.4cm,minimum height=1.15cm,
                  font=\footnotesize\sffamily},
  dblk/.style  ={blk,fill=decBg,draw=decPurple,minimum width=6.0cm},
  oblk/.style  ={blk,fill=outBg,draw=outGray,minimum width=3.8cm},
  op/.style     ={draw,rounded corners=2pt,minimum height=0.55cm,
                  font=\scriptsize\sffamily,align=center,thin,
                  inner sep=4pt,fill=white},
  dim/.style    ={font=\scriptsize\ttfamily,text=dimGray},
  frozen/.style ={font=\tiny\sffamily\itshape,text=frozenGray},
  stg/.style    ={font=\footnotesize\sffamily\bfseries,text=dimGray},
  arr/.style    ={-{Stealth[length=5pt,width=4pt]},thick,arrGray},
  sparr/.style  ={-{Stealth[length=5pt,width=4pt]},thick,sparseBlue,dashed},
  auxarr/.style ={-{Stealth[length=4pt,width=3pt]},semithick,lossRed,
                  densely dashed},
  gradarr/.style={-{Stealth[length=4pt,width=3pt]},semithick,gradGreen,
                  densely dotted},
}

\begin{tikzpicture}[node distance=0.8cm and 1.5cm]

% ═══════════════════════════════════════════════════════════════════════════════
%  Y-COORDINATE MAP (generous spacing)
%
%   0.0   Inputs
%  -1.8   Patch / Token Embedding
%  -3.4   Transformer Encoder backbones
%  -5.0   Linear projections
%  -6.8   Concatenation
%  -8.8   Pre-Norm Transformer (fusion)
% -10.8   Noisy Top-2 Router
% -13.2   Expert blocks
% -15.4   Weighted Aggregation
% -17.6   Decoder
% -19.4   Output Projection
% -20.4   Softmax
% -21.6   Answer
% ═══════════════════════════════════════════════════════════════════════════════


% ═══════════════════  STAGE 1 — INPUTS  ══════════════════════════════════════
\node[stg] at (-8.0,0) {\rotatebox{90}{INPUT}};

\node[inblk] (imgIn) at (-3.2,0) {\textbf{Image}~$I$};
\node[dim,below=4pt of imgIn] {$224{\times}224{\times}3$};

\node[inblk] (qIn) at (3.2,0) {\textbf{Question}~$Q$};
\node[dim,below=4pt of qIn] {$L$ tokens (Vietnamese)};


% ═══════════════════  STAGE 2/3 — ENCODERS  ═════════════════════════════════
\node[stg] at (-8.0,-3.4) {\rotatebox{90}{ENCODERS}};

% --- Visual column (left) ---
\node[vblk] (patchEmb) at (-3.2,-1.8)
  {Patch Embedding\\[-2pt]{\scriptsize $32{\times}32$ patches + \texttt{[CLS]}}};
\node[dim,right=8pt of patchEmb] {$50{\times}d_v$};

\node[vblk] (vitEnc) at (-3.2,-3.4)
  {\textbf{CLIP ViT-B/32}\\[-2pt]{\scriptsize Transformer Encoder}};
\node[frozen,right=8pt of vitEnc] {\textit{frozen}};

\node[op,draw=visBlue,fill=visBg!40] (visProj) at (-3.2,-5.0)
  {Linear~$W_v$};
\node[dim,right=8pt of visProj] {$50{\times}D$};

% --- Text column (right) ---
\node[tblk] (tokEmb) at (3.2,-1.8)
  {Token Embedding\\[-2pt]{\scriptsize + Position Encoding}};
\node[dim,left=8pt of tokEmb] {$L{\times}d_t$};

\node[tblk] (phoBERT) at (3.2,-3.4)
  {\textbf{PhoBERT-base}\\[-2pt]{\scriptsize Transformer Encoder}};
\node[frozen,left=8pt of phoBERT] {\textit{frozen}};

\node[op,draw=txtGreen,fill=txtBg!40] (txtProj) at (3.2,-5.0)
  {Linear~$W_t$};
\node[dim,left=8pt of txtProj] {$L{\times}D$};

% Encoder arrows
\draw[arr] (imgIn)    -- (patchEmb);
\draw[arr] (patchEmb) -- (vitEnc);
\draw[arr] (vitEnc)   -- (visProj);
\draw[arr] (qIn)      -- (tokEmb);
\draw[arr] (tokEmb)   -- (phoBERT);
\draw[arr] (phoBERT)  -- (txtProj);


% ═══════════════════  CONCATENATION  ═════════════════════════════════════════
\node[op,minimum width=5.2cm,draw=arrGray] (concat) at (0,-6.8)
  {Concatenation~$[\hat{V};\;\hat{T}]$};
\node[dim,below=4pt of concat] (concatDim) {$(N_v{+}L){\times}D$};

\draw[arr] (visProj.south) -- ++(0,-0.5) -| ([xshift=-1.2cm]concat.north);
\draw[arr] (txtProj.south) -- ++(0,-0.5) -| ([xshift= 1.2cm]concat.north);


% ═══════════════════  STAGE 4 — FUSION WITH MOE  ════════════════════════════
\node[stg] at (-8.0,-12.0) {\rotatebox{90}{FUSION~(MOE)}};

% (a) Pre-norm Transformer attention
\node[blk,fill=moeBg!60,draw=moeRed!70,minimum width=6.0cm]
  (preAttn) at (0,-8.8)
  {\textbf{Pre-Norm Transformer}~${\times}\,L_f{=}2$\\[-2pt]
   {\scriptsize MHSA~($H{=}8$)\;+\;FFN~($d_{\mathrm{ff}}{=}2048$)\;+\;LN}};
\node[dim,right=10pt of preAttn] {$(N_v{+}L){\times}D$};

% (b) Noisy Top-k Router
\node[rblk] (router) at (0,-10.8)
  {\textbf{Noisy Top-2 Router}\\[-2pt]
   {\scriptsize $\tilde{\ell}=W_r^\top x
     +\mathrm{Softplus}(W_n^\top x)\!\cdot\!\epsilon$\;,\;\;
     $\epsilon{\sim}\mathcal{N}(0,\sigma^2 I)$}\\[-2pt]
   {\scriptsize Softmax\;$\rightarrow$\;Top-2 selection per token}};
\node[dim,right=10pt of router,text=sparseBlue]
  {\textit{sparse:\;$k{=}2$ of $K{=}4$}};

% (c) Specialised Experts — wider horizontal spread
\node[eblk] (eVis)  at (-4.5,-13.2)
  {$E_{\mathrm{vis}}$\\[-1pt]\textbf{Vision}\\[-2pt]
   {\tiny Spatial MHSA}\\[-2pt]{\tiny + FFN}};
\node[eblk] (eTxt)  at (-1.5,-13.2)
  {$E_{\mathrm{txt}}$\\[-1pt]\textbf{Text}\\[-2pt]
   {\tiny Context MHSA}\\[-2pt]{\tiny + Masked FFN}};
\node[eblk] (eMM)   at ( 1.5,-13.2)
  {$E_{\mathrm{mm}}$\\[-1pt]\textbf{Multimodal}\\[-2pt]
   {\tiny Cross-Attn}\\[-2pt]{\tiny + Gate~$\gamma$}};
\node[eblk] (eSpec) at ( 4.5,-13.2)
  {$E_{\mathrm{spec}}$\\[-1pt]\textbf{Specialized}\\[-2pt]
   {\tiny Seg\,/\,Det\,/}\\[-2pt]{\tiny OCR\,/\,Scene}};

% (d) Weighted aggregation
\node[op,minimum width=7.0cm,draw=moeRed!70,fill=moeBg!50]
  (agg) at (0,-15.4)
  {Weighted Aggregation\;+\;LayerNorm};


% Load-balance loss (side branch — pushed further right)
\node[blk,fill=white,draw=lossRed,minimum width=2.6cm,
      font=\scriptsize\sffamily] (lbLoss) at (8.0,-10.8)
  {$\mathcal{L}_{\mathrm{lb}}
     =K\!\displaystyle\sum f_i\bar{p}_i$\\[0pt]
   {\tiny $\lambda_{\mathrm{lb}}{=}0.01$}};
\node[dim,above=5pt of lbLoss,text=lossRed]
  {\textit{Load-Balance Loss}};

% ── MOE highlighted background frame ────────────────────────────────────────
\begin{scope}[on background layer]
  \node[draw=moeBorder,fill=moeBg!20,rounded corners=6pt,
        very thick,inner sep=14pt,
        fit=(preAttn)(router)(eVis)(eSpec)(agg)(lbLoss)]
    (moeFrame) {};
  \node[font=\small\sffamily\bfseries,text=moeRed,
        anchor=north west]
    at ([xshift=6pt,yshift=-1pt]moeFrame.north west)
    {Cross-Modal Fusion with Mixture-of-Experts};
\end{scope}

% ── Arrows inside MOE ───────────────────────────────────────────────────────
\draw[arr] (concat) -- (preAttn);
\draw[arr] (preAttn) -- (router);

% Router → experts  (sparse dashed blue)
\draw[sparr] (router.south) -- ++(0,-0.5) -| (eVis.north);
\draw[sparr] (router.south) -- ++(0,-0.5) -| (eTxt.north);
\draw[sparr] (router.south) -- ++(0,-0.5) -| (eMM.north);
\draw[sparr] (router.south) -- ++(0,-0.5) -| (eSpec.north);

% Experts → aggregation
\draw[arr] (eVis.south)  -- ++(0,-0.45) -| ([xshift=-2.1cm]agg.north);
\draw[arr] (eTxt.south)  -- ++(0,-0.50) -| ([xshift=-0.7cm]agg.north);
\draw[arr] (eMM.south)   -- ++(0,-0.50) -| ([xshift= 0.7cm]agg.north);
\draw[arr] (eSpec.south) -- ++(0,-0.45) -| ([xshift= 2.1cm]agg.north);

% Router → load-balance loss
\draw[auxarr] (router.east) -- (lbLoss.west);


% ═══════════════════  STAGE 5 — DECODER  ════════════════════════════════════
\node[stg] at (-8.0,-18) {\rotatebox{90}{DECODER}};

\node[dblk] (decoder) at (0,-18)
  {\textbf{Autoregressive Transformer Decoder}~${\times}\,L_d{=}6$\\[-2pt]
   {\scriptsize Causal Self-Attn\;$\to$\;Cross-Attn\;$\to$\;FFN}\\[-2pt]
   {\scriptsize $H_d{=}8$ heads,\;\; $d_{\mathrm{ff}}{=}2048$,\;\; Pre-Norm}};
\node[dim,right=10pt of decoder]
  {\parbox{2.4cm}{\centering\scriptsize encoder memory\\[1pt]
   $(N_v{+}L){\times}D$}};

% Teacher-forcing input (left, with more clearance)
\node[op,draw=decPurple,fill=decBg!40,minimum width=2.0cm]
  (tfEmb) at (-5.8,-18)
  {\parbox{1.8cm}{\centering Embed.~$W_e$\\[1pt]{\tiny + Pos Enc}\\[1pt]
   {\tiny (teacher forcing)}}};
\node[dim,above=1pt of tfEmb] {$a_{<t}$};

\draw[arr] (agg) -- (decoder);
\draw[arr] (tfEmb) -- (decoder);


% ═══════════════════  STAGE 6 — OUTPUT  ═════════════════════════════════════
\node[stg] at (-8.0,-20.6) {\rotatebox{90}{OUTPUT}};

\node[op,draw=outGray,minimum width=4.5cm] (outProj) at (0,-19.4)
  {Output Projection~$W_e^\top$\;\;{\scriptsize (weight-tied)}};
\node[dim,right=10pt of outProj] {$T{\times}|\mathcal{V}|$};

\node[op,draw=outGray,minimum width=2.8cm] (softmax) at (0,-20.4)
  {Softmax};

\node[oblk] (answer) at (0,-21.6)
  {\textbf{Generated Answer}\\[-2pt]
   {\scriptsize Open-ended Vietnamese text}};

\draw[arr] (decoder) -- (outProj);
\draw[arr] (outProj) -- (softmax);
\draw[arr] (softmax) -- (answer);


% ═══════════════════  TRAINING LOSS  ═════════════════════════════════════════
\node[blk,fill=white,draw=lossRed,minimum width=3.4cm,
      font=\scriptsize\sffamily] (totalLoss) at (8.0,-18.0)
  {\textbf{Training Loss}\\[-2pt]
   $\mathcal{L}=\mathcal{L}_{\mathrm{ce}}
     +\lambda\,\mathcal{L}_{\mathrm{lb}}$\\[-2pt]
   {\tiny label smoothing~$\epsilon{=}0.1$}};

\draw[auxarr] (softmax.east) -| (totalLoss.south);
\draw[auxarr] (lbLoss.south) |- ([yshift=4pt]totalLoss.north);


% ═══════════════════  GRADIENT FLOW  ════════════════════════════════════════
\draw[gradarr,very thick]
  ([xshift=8.3cm]answer.east) -- ++(0,20.0)
  node[midway,right=4pt,font=\scriptsize\sffamily,
       text=gradGreen,rotate=90,anchor=south]
  {$\nabla$~Gradient flow (backpropagation)};


% ═══════════════════  FROZEN BACKBONE NOTE  ═════════════════════════════════
\node[font=\tiny\sffamily,text=frozenGray,align=center,
      fill=white,draw=frozenGray,rounded corners=2pt,inner sep=3pt]
  (frozenNote) at (-6.2,-3.4)
  {\textsf{\textbf{$\times$}}~Frozen backbones:\\no gradient update};

\draw[frozenGray,thin,densely dotted]
  (frozenNote.east) -- ++(0.2,0) |- (vitEnc.west);
\draw[frozenGray,thin,densely dotted]
  (frozenNote.east) -- ++(0.2,0) |- (phoBERT.west);


% ═══════════════════  LEGEND (top-right)  ════════════════════════════════════
\node[anchor=north west,font=\scriptsize\sffamily,text=dimGray,
      align=left,fill=white,draw=dimGray!50,
      rounded corners=3pt,inner sep=6pt]
  at (10.5, 0.5)
  {\textbf{Legend}\\[4pt]
   \tikz[baseline=-0.5ex]\draw[arr](0,0)--(0.6,0);
     ~Dense data flow\\[2pt]
   \tikz[baseline=-0.5ex]\draw[sparr](0,0)--(0.6,0);
     ~Sparse dispatch\\[2pt]
   \tikz[baseline=-0.5ex]\draw[auxarr](0,0)--(0.6,0);
     ~Auxiliary loss\\[2pt]
   \tikz[baseline=-0.5ex]\draw[gradarr](0,0)--(0.6,0);
     ~Gradient flow\\[5pt]
   \tikz\fill[visBg](0,0) rectangle (0.4,0.25);
     ~Visual pathway\\[2pt]
   \tikz\fill[txtBg](0,0) rectangle (0.4,0.25);
     ~Text pathway\\[2pt]
   \tikz\fill[moeBg](0,0) rectangle (0.4,0.25);
     ~MOE fusion (core)\\[2pt]
   \tikz\fill[decBg](0,0) rectangle (0.4,0.25);
     ~Decoder};


% ═══════════════════  COMPUTATIONAL PROPERTIES (right)  ═════════════════════
\node[anchor=north west,font=\scriptsize\sffamily,text=dimGray,
      align=left,fill=white,draw=dimGray!50,
      rounded corners=3pt,inner sep=6pt]
  at (10.5,-5.0)
  {\textbf{Computational Properties}\\[4pt]
   \textcolor{sparseBlue}{\rule{8pt}{2pt}}~%
     Sparse activation: Top-2 of $K{=}4$\\[2pt]
   \textcolor{sparseBlue}{\rule{8pt}{2pt}}~%
     Sparsity ratio: $k/K=50\%$ per token\\[2pt]
   Capacity factor: $\alpha=1.25$\\[2pt]
   $D=d_v=d_t=768$\\[2pt]
   $|\mathcal{V}|=64{,}000$ (PhoBERT vocab)\\[2pt]
   Weight tying: $W_e$ shared in\\
   \quad decoder embed.\ \& output proj.\\[4pt]
   \textcolor{gradGreen}{\rule[1pt]{8pt}{1.5pt}}~%
     Gradients flow through\\
   \quad fusion, MOE, decoder\\[2pt]
   \textcolor{frozenGray}{\textbf{$\times$}}~%
     No gradient to frozen encoders\\[4pt]
   \textit{Dense fusion:}\;$\mathcal{O}(S^2 D)$\\[1pt]
   \textit{Sparse MOE:}\;$\mathcal{O}(S\!\cdot\!k\!\cdot\!D)$};

\end{tikzpicture}
%
}
\caption{%
  Architecture of the proposed Generative VQA model with Mixture-of-Experts (MOE) fusion.
  Frozen CLIP ViT-B/32 and PhoBERT-base encode the image and Vietnamese question, respectively.
  Their projected embeddings are concatenated and processed through two pre-norm transformer layers,
  then routed via a Noisy Top-2 gating mechanism to $K{=}4$ heterogeneous specialised experts
  (Vision, Text, Multimodal, Specialized).
  Only $k{=}2$ experts are activated per token (50\% sparsity),
  maintaining computational efficiency while enabling adaptive multi-modal reasoning.
  The weighted expert outputs are decoded by a 6-layer autoregressive transformer with weight-tied embeddings.
  Dashed blue arrows denote sparse dispatch; dashed red arrows denote auxiliary loss paths.
  Gradients flow through the fusion, MOE, and decoder modules (dotted green);
  backbone encoders remain frozen.%
}
\label{fig:architecture}
\end{figure*}

% Place inside your paper .tex file — requires tikz + xcolor already loaded
% Packages needed in preamble:
%   \usepackage{tikz}
%   \usetikzlibrary{arrows.meta,positioning,shapes.geometric,calc,
%                    fit,backgrounds,decorations.pathreplacing,patterns}
% =============================================================================

\begin{figure*}[t]
\centering
\resizebox{\textwidth}{!}{%
% =============================================================================
% Architecture Diagram Body  (no preamble — for \input inside a figure)
% Requires in main document preamble:
%   \usepackage{tikz,amsmath,amssymb,xcolor}
%   \usetikzlibrary{arrows.meta,positioning,shapes.geometric,calc,
%                    fit,backgrounds,decorations.pathreplacing,patterns}
% =============================================================================

% ── Colours ──────────────────────────────────────────────────────────────────
\definecolor{visBlue}{HTML}{2E6DAD}
\definecolor{visBg}{HTML}{D6E6F5}
\definecolor{txtGreen}{HTML}{237A45}
\definecolor{txtBg}{HTML}{D4EDDE}
\definecolor{moeRed}{HTML}{C0392B}
\definecolor{moeBg}{HTML}{FDECEC}
\definecolor{moeBorder}{HTML}{D45A5A}
\definecolor{routerOrange}{HTML}{D98B2B}
\definecolor{routerBg}{HTML}{FDF3E2}
\definecolor{expPurple}{HTML}{7B5EA7}
\definecolor{expBg}{HTML}{EDE6F6}
\definecolor{decPurple}{HTML}{6A4C93}
\definecolor{decBg}{HTML}{E8E0F0}
\definecolor{outGray}{HTML}{4A4A4A}
\definecolor{outBg}{HTML}{F0F0F0}
\definecolor{arrGray}{HTML}{555555}
\definecolor{dimGray}{HTML}{888888}
\definecolor{lossRed}{HTML}{B03A2E}
\definecolor{gradGreen}{HTML}{2E7D32}
\definecolor{sparseBlue}{HTML}{1565C0}
\definecolor{frozenGray}{HTML}{95A5A6}
\definecolor{inputBlue}{HTML}{4A90D9}
\definecolor{inputBg}{HTML}{EBF2FA}

% ── Styles ───────────────────────────────────────────────────────────────────
\tikzset{
  blk/.style={draw,rounded corners=3pt,minimum height=0.95cm,
              font=\small\sffamily,align=center,thick,inner sep=6pt},
  inblk/.style ={blk,fill=inputBg,draw=inputBlue,minimum width=3.2cm},
  vblk/.style  ={blk,fill=visBg,draw=visBlue,minimum width=3.6cm},
  tblk/.style  ={blk,fill=txtBg,draw=txtGreen,minimum width=3.6cm},
  rblk/.style  ={blk,fill=routerBg,draw=routerOrange,minimum width=5.0cm},
  eblk/.style  ={blk,fill=expBg,draw=expPurple,
                  minimum width=2.4cm,minimum height=1.15cm,
                  font=\footnotesize\sffamily},
  dblk/.style  ={blk,fill=decBg,draw=decPurple,minimum width=6.0cm},
  oblk/.style  ={blk,fill=outBg,draw=outGray,minimum width=3.8cm},
  op/.style     ={draw,rounded corners=2pt,minimum height=0.55cm,
                  font=\scriptsize\sffamily,align=center,thin,
                  inner sep=4pt,fill=white},
  dim/.style    ={font=\scriptsize\ttfamily,text=dimGray},
  frozen/.style ={font=\tiny\sffamily\itshape,text=frozenGray},
  stg/.style    ={font=\footnotesize\sffamily\bfseries,text=dimGray},
  arr/.style    ={-{Stealth[length=5pt,width=4pt]},thick,arrGray},
  sparr/.style  ={-{Stealth[length=5pt,width=4pt]},thick,sparseBlue,dashed},
  auxarr/.style ={-{Stealth[length=4pt,width=3pt]},semithick,lossRed,
                  densely dashed},
  gradarr/.style={-{Stealth[length=4pt,width=3pt]},semithick,gradGreen,
                  densely dotted},
}

\begin{tikzpicture}[node distance=0.8cm and 1.5cm]

% ═══════════════════════════════════════════════════════════════════════════════
%  Y-COORDINATE MAP (generous spacing)
%
%   0.0   Inputs
%  -1.8   Patch / Token Embedding
%  -3.4   Transformer Encoder backbones
%  -5.0   Linear projections
%  -6.8   Concatenation
%  -8.8   Pre-Norm Transformer (fusion)
% -10.8   Noisy Top-2 Router
% -13.2   Expert blocks
% -15.4   Weighted Aggregation
% -17.6   Decoder
% -19.4   Output Projection
% -20.4   Softmax
% -21.6   Answer
% ═══════════════════════════════════════════════════════════════════════════════


% ═══════════════════  STAGE 1 — INPUTS  ══════════════════════════════════════
\node[stg] at (-8.0,0) {\rotatebox{90}{INPUT}};

\node[inblk] (imgIn) at (-3.2,0) {\textbf{Image}~$I$};
\node[dim,below=4pt of imgIn] {$224{\times}224{\times}3$};

\node[inblk] (qIn) at (3.2,0) {\textbf{Question}~$Q$};
\node[dim,below=4pt of qIn] {$L$ tokens (Vietnamese)};


% ═══════════════════  STAGE 2/3 — ENCODERS  ═════════════════════════════════
\node[stg] at (-8.0,-3.4) {\rotatebox{90}{ENCODERS}};

% --- Visual column (left) ---
\node[vblk] (patchEmb) at (-3.2,-1.8)
  {Patch Embedding\\[-2pt]{\scriptsize $32{\times}32$ patches + \texttt{[CLS]}}};
\node[dim,right=8pt of patchEmb] {$50{\times}d_v$};

\node[vblk] (vitEnc) at (-3.2,-3.4)
  {\textbf{CLIP ViT-B/32}\\[-2pt]{\scriptsize Transformer Encoder}};
\node[frozen,right=8pt of vitEnc] {\textit{frozen}};

\node[op,draw=visBlue,fill=visBg!40] (visProj) at (-3.2,-5.0)
  {Linear~$W_v$};
\node[dim,right=8pt of visProj] {$50{\times}D$};

% --- Text column (right) ---
\node[tblk] (tokEmb) at (3.2,-1.8)
  {Token Embedding\\[-2pt]{\scriptsize + Position Encoding}};
\node[dim,left=8pt of tokEmb] {$L{\times}d_t$};

\node[tblk] (phoBERT) at (3.2,-3.4)
  {\textbf{PhoBERT-base}\\[-2pt]{\scriptsize Transformer Encoder}};
\node[frozen,left=8pt of phoBERT] {\textit{frozen}};

\node[op,draw=txtGreen,fill=txtBg!40] (txtProj) at (3.2,-5.0)
  {Linear~$W_t$};
\node[dim,left=8pt of txtProj] {$L{\times}D$};

% Encoder arrows
\draw[arr] (imgIn)    -- (patchEmb);
\draw[arr] (patchEmb) -- (vitEnc);
\draw[arr] (vitEnc)   -- (visProj);
\draw[arr] (qIn)      -- (tokEmb);
\draw[arr] (tokEmb)   -- (phoBERT);
\draw[arr] (phoBERT)  -- (txtProj);


% ═══════════════════  CONCATENATION  ═════════════════════════════════════════
\node[op,minimum width=5.2cm,draw=arrGray] (concat) at (0,-6.8)
  {Concatenation~$[\hat{V};\;\hat{T}]$};
\node[dim,below=4pt of concat] (concatDim) {$(N_v{+}L){\times}D$};

\draw[arr] (visProj.south) -- ++(0,-0.5) -| ([xshift=-1.2cm]concat.north);
\draw[arr] (txtProj.south) -- ++(0,-0.5) -| ([xshift= 1.2cm]concat.north);


% ═══════════════════  STAGE 4 — FUSION WITH MOE  ════════════════════════════
\node[stg] at (-8.0,-12.0) {\rotatebox{90}{FUSION~(MOE)}};

% (a) Pre-norm Transformer attention
\node[blk,fill=moeBg!60,draw=moeRed!70,minimum width=6.0cm]
  (preAttn) at (0,-8.8)
  {\textbf{Pre-Norm Transformer}~${\times}\,L_f{=}2$\\[-2pt]
   {\scriptsize MHSA~($H{=}8$)\;+\;FFN~($d_{\mathrm{ff}}{=}2048$)\;+\;LN}};
\node[dim,right=10pt of preAttn] {$(N_v{+}L){\times}D$};

% (b) Noisy Top-k Router
\node[rblk] (router) at (0,-10.8)
  {\textbf{Noisy Top-2 Router}\\[-2pt]
   {\scriptsize $\tilde{\ell}=W_r^\top x
     +\mathrm{Softplus}(W_n^\top x)\!\cdot\!\epsilon$\;,\;\;
     $\epsilon{\sim}\mathcal{N}(0,\sigma^2 I)$}\\[-2pt]
   {\scriptsize Softmax\;$\rightarrow$\;Top-2 selection per token}};
\node[dim,right=10pt of router,text=sparseBlue]
  {\textit{sparse:\;$k{=}2$ of $K{=}4$}};

% (c) Specialised Experts — wider horizontal spread
\node[eblk] (eVis)  at (-4.5,-13.2)
  {$E_{\mathrm{vis}}$\\[-1pt]\textbf{Vision}\\[-2pt]
   {\tiny Spatial MHSA}\\[-2pt]{\tiny + FFN}};
\node[eblk] (eTxt)  at (-1.5,-13.2)
  {$E_{\mathrm{txt}}$\\[-1pt]\textbf{Text}\\[-2pt]
   {\tiny Context MHSA}\\[-2pt]{\tiny + Masked FFN}};
\node[eblk] (eMM)   at ( 1.5,-13.2)
  {$E_{\mathrm{mm}}$\\[-1pt]\textbf{Multimodal}\\[-2pt]
   {\tiny Cross-Attn}\\[-2pt]{\tiny + Gate~$\gamma$}};
\node[eblk] (eSpec) at ( 4.5,-13.2)
  {$E_{\mathrm{spec}}$\\[-1pt]\textbf{Specialized}\\[-2pt]
   {\tiny Seg\,/\,Det\,/}\\[-2pt]{\tiny OCR\,/\,Scene}};

% (d) Weighted aggregation
\node[op,minimum width=7.0cm,draw=moeRed!70,fill=moeBg!50]
  (agg) at (0,-15.4)
  {Weighted Aggregation\;+\;LayerNorm};


% Load-balance loss (side branch — pushed further right)
\node[blk,fill=white,draw=lossRed,minimum width=2.6cm,
      font=\scriptsize\sffamily] (lbLoss) at (8.0,-10.8)
  {$\mathcal{L}_{\mathrm{lb}}
     =K\!\displaystyle\sum f_i\bar{p}_i$\\[0pt]
   {\tiny $\lambda_{\mathrm{lb}}{=}0.01$}};
\node[dim,above=5pt of lbLoss,text=lossRed]
  {\textit{Load-Balance Loss}};

% ── MOE highlighted background frame ────────────────────────────────────────
\begin{scope}[on background layer]
  \node[draw=moeBorder,fill=moeBg!20,rounded corners=6pt,
        very thick,inner sep=14pt,
        fit=(preAttn)(router)(eVis)(eSpec)(agg)(lbLoss)]
    (moeFrame) {};
  \node[font=\small\sffamily\bfseries,text=moeRed,
        anchor=north west]
    at ([xshift=6pt,yshift=-1pt]moeFrame.north west)
    {Cross-Modal Fusion with Mixture-of-Experts};
\end{scope}

% ── Arrows inside MOE ───────────────────────────────────────────────────────
\draw[arr] (concat) -- (preAttn);
\draw[arr] (preAttn) -- (router);

% Router → experts  (sparse dashed blue)
\draw[sparr] (router.south) -- ++(0,-0.5) -| (eVis.north);
\draw[sparr] (router.south) -- ++(0,-0.5) -| (eTxt.north);
\draw[sparr] (router.south) -- ++(0,-0.5) -| (eMM.north);
\draw[sparr] (router.south) -- ++(0,-0.5) -| (eSpec.north);

% Experts → aggregation
\draw[arr] (eVis.south)  -- ++(0,-0.45) -| ([xshift=-2.1cm]agg.north);
\draw[arr] (eTxt.south)  -- ++(0,-0.50) -| ([xshift=-0.7cm]agg.north);
\draw[arr] (eMM.south)   -- ++(0,-0.50) -| ([xshift= 0.7cm]agg.north);
\draw[arr] (eSpec.south) -- ++(0,-0.45) -| ([xshift= 2.1cm]agg.north);

% Router → load-balance loss
\draw[auxarr] (router.east) -- (lbLoss.west);


% ═══════════════════  STAGE 5 — DECODER  ════════════════════════════════════
\node[stg] at (-8.0,-18) {\rotatebox{90}{DECODER}};

\node[dblk] (decoder) at (0,-18)
  {\textbf{Autoregressive Transformer Decoder}~${\times}\,L_d{=}6$\\[-2pt]
   {\scriptsize Causal Self-Attn\;$\to$\;Cross-Attn\;$\to$\;FFN}\\[-2pt]
   {\scriptsize $H_d{=}8$ heads,\;\; $d_{\mathrm{ff}}{=}2048$,\;\; Pre-Norm}};
\node[dim,right=10pt of decoder]
  {\parbox{2.4cm}{\centering\scriptsize encoder memory\\[1pt]
   $(N_v{+}L){\times}D$}};

% Teacher-forcing input (left, with more clearance)
\node[op,draw=decPurple,fill=decBg!40,minimum width=2.0cm]
  (tfEmb) at (-5.8,-18)
  {\parbox{1.8cm}{\centering Embed.~$W_e$\\[1pt]{\tiny + Pos Enc}\\[1pt]
   {\tiny (teacher forcing)}}};
\node[dim,above=1pt of tfEmb] {$a_{<t}$};

\draw[arr] (agg) -- (decoder);
\draw[arr] (tfEmb) -- (decoder);


% ═══════════════════  STAGE 6 — OUTPUT  ═════════════════════════════════════
\node[stg] at (-8.0,-20.6) {\rotatebox{90}{OUTPUT}};

\node[op,draw=outGray,minimum width=4.5cm] (outProj) at (0,-19.4)
  {Output Projection~$W_e^\top$\;\;{\scriptsize (weight-tied)}};
\node[dim,right=10pt of outProj] {$T{\times}|\mathcal{V}|$};

\node[op,draw=outGray,minimum width=2.8cm] (softmax) at (0,-20.4)
  {Softmax};

\node[oblk] (answer) at (0,-21.6)
  {\textbf{Generated Answer}\\[-2pt]
   {\scriptsize Open-ended Vietnamese text}};

\draw[arr] (decoder) -- (outProj);
\draw[arr] (outProj) -- (softmax);
\draw[arr] (softmax) -- (answer);


% ═══════════════════  TRAINING LOSS  ═════════════════════════════════════════
\node[blk,fill=white,draw=lossRed,minimum width=3.4cm,
      font=\scriptsize\sffamily] (totalLoss) at (8.0,-18.0)
  {\textbf{Training Loss}\\[-2pt]
   $\mathcal{L}=\mathcal{L}_{\mathrm{ce}}
     +\lambda\,\mathcal{L}_{\mathrm{lb}}$\\[-2pt]
   {\tiny label smoothing~$\epsilon{=}0.1$}};

\draw[auxarr] (softmax.east) -| (totalLoss.south);
\draw[auxarr] (lbLoss.south) |- ([yshift=4pt]totalLoss.north);


% ═══════════════════  GRADIENT FLOW  ════════════════════════════════════════
\draw[gradarr,very thick]
  ([xshift=8.3cm]answer.east) -- ++(0,20.0)
  node[midway,right=4pt,font=\scriptsize\sffamily,
       text=gradGreen,rotate=90,anchor=south]
  {$\nabla$~Gradient flow (backpropagation)};


% ═══════════════════  FROZEN BACKBONE NOTE  ═════════════════════════════════
\node[font=\tiny\sffamily,text=frozenGray,align=center,
      fill=white,draw=frozenGray,rounded corners=2pt,inner sep=3pt]
  (frozenNote) at (-6.2,-3.4)
  {\textsf{\textbf{$\times$}}~Frozen backbones:\\no gradient update};

\draw[frozenGray,thin,densely dotted]
  (frozenNote.east) -- ++(0.2,0) |- (vitEnc.west);
\draw[frozenGray,thin,densely dotted]
  (frozenNote.east) -- ++(0.2,0) |- (phoBERT.west);


% ═══════════════════  LEGEND (top-right)  ════════════════════════════════════
\node[anchor=north west,font=\scriptsize\sffamily,text=dimGray,
      align=left,fill=white,draw=dimGray!50,
      rounded corners=3pt,inner sep=6pt]
  at (10.5, 0.5)
  {\textbf{Legend}\\[4pt]
   \tikz[baseline=-0.5ex]\draw[arr](0,0)--(0.6,0);
     ~Dense data flow\\[2pt]
   \tikz[baseline=-0.5ex]\draw[sparr](0,0)--(0.6,0);
     ~Sparse dispatch\\[2pt]
   \tikz[baseline=-0.5ex]\draw[auxarr](0,0)--(0.6,0);
     ~Auxiliary loss\\[2pt]
   \tikz[baseline=-0.5ex]\draw[gradarr](0,0)--(0.6,0);
     ~Gradient flow\\[5pt]
   \tikz\fill[visBg](0,0) rectangle (0.4,0.25);
     ~Visual pathway\\[2pt]
   \tikz\fill[txtBg](0,0) rectangle (0.4,0.25);
     ~Text pathway\\[2pt]
   \tikz\fill[moeBg](0,0) rectangle (0.4,0.25);
     ~MOE fusion (core)\\[2pt]
   \tikz\fill[decBg](0,0) rectangle (0.4,0.25);
     ~Decoder};


% ═══════════════════  COMPUTATIONAL PROPERTIES (right)  ═════════════════════
\node[anchor=north west,font=\scriptsize\sffamily,text=dimGray,
      align=left,fill=white,draw=dimGray!50,
      rounded corners=3pt,inner sep=6pt]
  at (10.5,-5.0)
  {\textbf{Computational Properties}\\[4pt]
   \textcolor{sparseBlue}{\rule{8pt}{2pt}}~%
     Sparse activation: Top-2 of $K{=}4$\\[2pt]
   \textcolor{sparseBlue}{\rule{8pt}{2pt}}~%
     Sparsity ratio: $k/K=50\%$ per token\\[2pt]
   Capacity factor: $\alpha=1.25$\\[2pt]
   $D=d_v=d_t=768$\\[2pt]
   $|\mathcal{V}|=64{,}000$ (PhoBERT vocab)\\[2pt]
   Weight tying: $W_e$ shared in\\
   \quad decoder embed.\ \& output proj.\\[4pt]
   \textcolor{gradGreen}{\rule[1pt]{8pt}{1.5pt}}~%
     Gradients flow through\\
   \quad fusion, MOE, decoder\\[2pt]
   \textcolor{frozenGray}{\textbf{$\times$}}~%
     No gradient to frozen encoders\\[4pt]
   \textit{Dense fusion:}\;$\mathcal{O}(S^2 D)$\\[1pt]
   \textit{Sparse MOE:}\;$\mathcal{O}(S\!\cdot\!k\!\cdot\!D)$};

\end{tikzpicture}
%
}
\caption{%
  Architecture of the proposed Generative VQA model with Mixture-of-Experts (MOE) fusion.
  Frozen CLIP ViT-B/32 and PhoBERT-base encode the image and Vietnamese question, respectively.
  Their projected embeddings are concatenated and processed through two pre-norm transformer layers,
  then routed via a Noisy Top-2 gating mechanism to $K{=}4$ heterogeneous specialised experts
  (Vision, Text, Multimodal, Specialized).
  Only $k{=}2$ experts are activated per token (50\% sparsity),
  maintaining computational efficiency while enabling adaptive multi-modal reasoning.
  The weighted expert outputs are decoded by a 6-layer autoregressive transformer with weight-tied embeddings.
  Dashed blue arrows denote sparse dispatch; dashed red arrows denote auxiliary loss paths.
  Gradients flow through the fusion, MOE, and decoder modules (dotted green);
  backbone encoders remain frozen.%
}
\label{fig:architecture}
\end{figure*}
